\section{Introdução}
\label{sec:introducao}

A administração pública brasileira exerce papel central na prestação de serviços essenciais à população. É no âmbito municipal que se concretizam as políticas públicas cotidianas, abrangendo desde a atenção primária à saúde e a educação básica até a manutenção da malha viária urbana e os serviços de saneamento. A execução contínua dessas atribuições demanda gestão eficiente e mobilização de expressivos recursos públicos, viabilizados tanto por receitas tributárias próprias quanto por transferências constitucionais das esferas estadual e federal.

Para atender às demandas da coletividade, o poder público recorre rotineiramente ao mercado para aquisição de insumos médico-hospitalares, contratação de obras de engenharia civil, locação de frotas e prestação de serviços especializados. Dessa forma, consolida-se a interface entre o ente público contratante (prefeituras e secretarias municipais) e os fornecedores privados. Embora a legislação nacional — estabelecida pela Nova Lei de Licitações e Contratos Administrativos (Lei Federal nº 14.133/2021)~\cite{brasil2021licitacoes}, pela Lei de Defesa da Concorrência (Lei Federal nº 12.529/2011)~\cite{brasil2011concorrencia} e pela Lei Anticorrupção (Lei Federal nº 12.846/2013)~\cite{brasil2013anticorrupcao} — imponha rígidos princípios de publicidade, isonomia e ampla competitividade, práticas anticompetitivas e arranjos societários ou comerciais podem explorar brechas relacionais e temporais, dificultando sua detecção pelos mecanismos tradicionais de controle.

Essa dificuldade decorre do fato de que instrumentos convencionais de fiscalização operam sob perspectiva predominantemente tabular, avaliando certames, contratos e empenhos de maneira atômica e isolada. Sob essa visão fragmentada, verificações lineares têm baixa capacidade de capturar padrões sistêmicos complexos, tais como vínculos societários ou cadastrais ocultos entre empresas formalmente concorrentes, o rodízio coordenado de fornecedores vencedores com propostas de cobertura (\textit{cover bidding}) e a concentração orçamentária em poucos credores ou grupos econômicos. O sinal investigativo relevante, portanto, emerge do comportamento cumulativo e das relações estruturais consolidadas entre o poder público e o mercado fornecedor ao longo do tempo.

Para superar essa limitação, a Teoria dos Grafos fornece uma base analítica rigorosa. Ao modelar unidades orçamentárias, empresas licitantes e seus sócios reais como vértices ($V$) e ao representar participações, co-disputas, alternâncias de vitórias e fluxos empenhados como arestas ponderadas ($E$), a abordagem em redes transcende a inspeção isolada de tabelas. Essa perspectiva possibilita mapear conexões indiretas, isolar subgrupos suspeitos de co-atuação e quantificar assimetrias de dependência e concentração financeira que permaneceriam ocultas em auditorias convencionais.

A motivação prática para este estudo nasce da lacuna de granularidade observada nos portais locais de transparência: enquanto o portal municipal de Divinópolis/MG disponibiliza apenas visões agregadas~\cite{portal_transparencia_betha}, os dados jurisdicionados do Sistema Informatizado de Contas dos Municípios (SICOM), mantidos em dados abertos pelo Tribunal de Contas do Estado de Minas Gerais (TCE-MG)~\cite{tce_mg_dados_abertos,tcemg2025in01}, reúnem 55 tabelas relacionais que cobrem todo o ciclo da despesa pública. Assim, este trabalho tem como objetivo geral desenvolver e validar um \textit{pipeline} computacional baseado em Teoria dos Grafos para \textbf{triagem indiciária de risco} em licitações e contratos administrativos do município de Divinópolis (MG) no exercício consolidado de 2025. Em consonância com a jurisprudência do Tribunal de Contas da União (TCU)~\cite{tcu2011acordao2341,tcu2024acordao1798}, enfatiza-se que os indicadores estruturais obtidos não constituem prova automática de fraude ou ilícito, mas sim critérios quantitativos para priorização de auditoria.

De forma específica, a investigação estrutura-se em três eixos complementares: (i) no \textbf{Eixo Societário e de Materialidade}, construir um grafo heterogêneo multicamada cruzando co-disputas com o Quadro de Sócios e Administradores (QSA) da Receita Federal para pontuar vínculos ocultos, isolar triângulos societários ($K_3$) e ponderar o risco pelo valor efetivamente empenhado e pago; (ii) no \textbf{Eixo Comportamental}, modelar um dígrafo de disputas (Perdedor $\to$ Vencedor) para identificar inversões recíprocas de vitórias em datas distintas, ciclos curtos ($C_3/C_4$), comunidades de co-concorrência e supressão de desconto frente às cotações de referência; e (iii) no \textbf{Eixo de Concentração Orçamentária}, estruturar um grafo bipartido Unidade Orçamentária $\leftrightarrow$ Credor para mensurar o Índice Herfindahl--Hirschman (HHI), a entropia de dispersão de gastos, o \textit{PageRank} ponderado e o impacto de grupos societários ($\Delta\text{HHI}$), diferenciando a dependência estrutural de entes públicos ou filantrópicos da concentração em fornecedores privados.

O restante deste artigo está organizado da seguinte maneira. A Seção~\ref{sec:referencial} apresenta o referencial teórico, integrando o arcabouço jurídico-regulatório, as tipologias de risco em compras públicas e os algoritmos clássicos de grafos empregados. A Seção~\ref{sec:metodologia} detalha a arquitetura do \textit{pipeline} em dados do SICOM 2025, a modelagem formal de vértices e arestas, o protocolo de conformidade com a LGPD e o fluxograma completo da solução. A Seção~\ref{sec:resultados} expõe e discute criticamente os resultados empíricos dos três eixos, os achados consolidados, os gargalos computacionais e os tempos reais de execução. Por fim, a Seção~\ref{sec:conclusao} sintetiza as contribuições obtidas, analisa as limitações empíricas da base pública e aponta direções para trabalhos futuros. Com esta proposta, espera-se disponibilizar uma metodologia reprodutível, auditável e escalável capaz de apoiar órgãos de controle interno e externo na triagem preventiva de anomalias contratuais no setor público.
